\section{Full integral homology}
Let \(G=(V,E)\) be finite, simple and connected, with at least one edge.
At \(v\), take \(\Sigma_{0,d_v+2}\times S^1\). Cap the two designated
ports by the primitive meridians \(p_vz+h\) and \(-p_vz+h\), where
\(p_v\) is a positive integer. Glue the remaining ports by
\(z_{ve}=\epsilon_eh_w,\ h_v=\epsilon_ez_{we}\), with
\(\epsilon_e\in\{\pm1\}\). Denote the closed oriented manifold by
\(M_{G,\epsilon,\mathbf p}\). Let \(A_\epsilon\) be the signed
adjacency matrix, \(P=\diag(p_v)\), and \(b_G=|E|-|V|+1\).

\begin{theorem}[Integral formula]
There is a noncanonical isomorphism
\[
 H_1(M_{G,\epsilon,\mathbf p};\mathbb Z)
 \cong \mathbb Z^{b_G}\oplus\coker_{\mathbb Z}(P A_\epsilon P).
\]
In particular, the paired \(+2,-2\) family of Paper 50B satisfies
\[
 H_1(M_{G,\epsilon};\mathbb Z)
 \cong \mathbb Z^{b_G}\oplus\coker_{\mathbb Z}(4A_\epsilon).
\]
\end{theorem}
\begin{proof}
Use generators \(a_v,b_v\) for the two cap boundary sections, \(h_v\)
for the fiber and \(z_{ve}\) for the remaining sections. The block and
cap relations are
\[
 p_va_v+h_v=0,\qquad -p_vb_v+h_v=0,\qquad
 a_v+b_v+\sum_{e\ni v}z_{ve}=0.
\]
The coefficients of \(h_v\) in the first relation and of \(b_v\) in
the last are units over \(\mathbb Z\), so eliminate
\(h_v=-p_va_v\) and \(b_v=-a_v-\sum_ez_{ve}\) integrally.
The remaining cap relation is exactly
\(p_v\sum_ez_{ve}=0\). Each seam gives
\(z_{ve}=-\epsilon_ep_wa_w\); eliminating these generators leaves
\[
 p_v\sum_w\epsilon_{vw}p_wa_w=0
\]
at every vertex. This presents \(\coker(PA_\epsilon P)\).
The graph-of-spaces Mayer--Vietoris sequence has the remaining term
\(\ker(\mathbb Z^E\to\mathbb Z^V)=\mathbb Z^{b_G}\), because
seams identify point classes with degree one regardless of their sign.
Thus there is a short exact sequence with that free quotient; it splits.
No division by \(2\), or by any \(p_v\), occurs.
\end{proof}

\begin{corollary}[Smith form for complete bipartite backbones]
Let \(G=K_{N,N}\), let \(D\) be a nonsingular binary
\((N-1)\times(N-1)\) matrix, and set
\[
 B(D)=\begin{pmatrix}1&\mathbf1^T\\
 \mathbf1&\mathbf1\mathbf1^T-2D\end{pmatrix}.
\]
If the positive Smith factors of \(D\) are \(d_1\mid\cdots\mid d_{N-1}\),
then for the \(+2,-2\) caps
\[
 H_1(M_{B(D)};\mathbb Z)
 \cong \mathbb Z^{(N-1)^2}\oplus(\mathbb Z/4)^2
 \oplus\bigoplus_{i=1}^{N-1}(\mathbb Z/(8d_i))^2.
\]
\end{corollary}
\begin{proof}
Integral row and column subtraction gives
\(B(D)\sim\diag(1,-2D)\), so its Smith factors are
\(1,2d_1,\ldots,2d_{N-1}\). For a bipartite adjacency
\(A=\left(\begin{smallmatrix}0&B\\B^T&0\end{smallmatrix}\right)\),
a column permutation gives \(B\oplus B^T\); hence its Smith factors
are those of \(B\) doubled. Multiplication by \(4\) multiplies
every nonzero Smith factor by \(4\). Finally
\(b_G=N^2-2N+1\).
\end{proof}

\section{What the torus kernel makes invisible}
\begin{proposition}[Variable caps and local ribbon insertions]
Fix the positive cap orders \(p_v\). If \(G\) is bipartite, then the
oriented \(\mathrm{SU}(2)\) RT invariant is independent of every seam
sign \(\epsilon_e\), at every admissible level. The same holds with
arbitrary fixed colored framed ribbon graphs placed in the interiors of
the capped blocks, provided that no ribbon intersects a sewing torus.
More generally the assertion holds for any modular category whose simple
objects are all self-dual, in the same extended-TQFT normalization.
\end{proposition}
\begin{proof}
Each block with any such interior decoration determines a tensor in the
tensor product of its unpunctured torus state spaces. Reversing a seam sign
inserts the torus involution \(-I\). With the section Lagrangian and its
zero-weight lift, this operator is charge conjugation \(C=S^2\), hence
the identity when all simple objects are self-dual. Thus the bare tensor
contraction is identical for all signings, regardless of the individual
block tensors.

For completeness the oriented anomaly also remains unchanged. In the
port notation of Paper 50B, the cap/seam matrix \(T\) has cap entries
\(p_v,-p_v\) and seam blocks
\(\left(\begin{smallmatrix}0&-\epsilon_e\\-\epsilon_e&0\end{smallmatrix}\right)\).
Consequently \(\sig T=0\) and
\(UT^{-1}U^T=-A_\epsilon\), since the two reciprocal cap entries
cancel at every vertex. The Maslov correction is therefore
\(\sig A_\epsilon\), exactly as for cap orders two. A bipartite
diagonal involution conjugates \(A_\epsilon\) to \(-A_\epsilon\),
so this signature vanishes. Interior ribbon decorations do not change
the homological Lagrangians or their Maslov index.
\end{proof}
\begin{remark}
This does not concern arbitrary links crossing the sewing tori. Such a
cut has punctured boundary state spaces, on which the preceding
unpunctured-torus identity is not the relevant assertion. In particular,
this proposition cannot be used to declare all possible colored
observables blind.
\end{remark}

\section{All odd pointed theories: the phase is fixed as well}
Return to the cap orders \(p_v=2\). Let \((\mathcal A,Q)\) be a
nondegenerate finite metric group of odd order \(m\), and use its unitary
pointed modular category. Write
\[
 \beta(x,y)=\frac{Q(x+y)}{Q(x)Q(y)},\qquad
 S_{xy}=m^{-1/2}\beta(x,y),\qquad \Theta_{xx}=Q(x).
\]
Changing the convention to the complex conjugate \(S\) makes no
difference to the result. Let \(\mathcal P_{\mathcal A,Q}\) be the
vacuum-divided invariant, with \(\mathcal P(S^3)=1\).

\begin{theorem}[Positive odd-pointed scalar]
If \(G\) is bipartite with parts of sizes \(n_A,n_B\) and signed
biadjacency \(B\), then
\[
 \mathcal P_{\mathcal A,Q}(M_{G,\epsilon})
 =m^{(|E|+1)/2-n_B}\,
 \left|\ker\bigl(B:\mathcal A^{n_B}\longrightarrow
                         \mathcal A^{n_A}\bigr)\right|
 =\sqrt{|\Hom(H_1(M_{G,\epsilon};\mathbb Z),\mathcal A)|}.
\]
In particular this is a positive real number and depends on the group
\(\mathcal A\), but not on its nondegenerate quadratic form \(Q\).
\end{theorem}
\begin{proof}
The positive cap amplitude at color \(a\) is
\[
 c_+(a)=m^{-1}\sum_{x\in\mathcal A}\beta(a,x)Q(x)^2.
\]
The negative cap amplitude is its complex conjugate: use
\(Q(-x)=Q(x)\) and change \(x\) to \(-x\). Expanding the squared
modulus and putting \(x=y+c\) gives
\[
 |c_+(a)|^2=m^{-2}\sum_c\beta(a,c)Q(c)^2
                    \sum_y\beta(y,c)^2=\frac1m.
\]
Indeed nondegeneracy says the inner character sum vanishes unless
\(2c=0\), and odd order then forces \(c=0\). Thus every paired cap
has the same positive weight \(m^{-1}\).

Since \(S_{0a}=m^{-1/2}\), a degree-\(d_v\) vertex contributes
\(m^{d_v/2-1}\). A seam contributes
\(m^{-1/2}\beta(x_v,x_w)^{\epsilon_e}\), with simultaneous duals
if the opposite orientation convention is chosen. The oriented Maslov
phase is one by bipartiteness. Dividing by \(Z(S^3)=m^{-1/2}\)
therefore gives
\[
 \mathcal P=m^{(|E|+1)/2-|V|}
 \sum_{x\in\mathcal A^{n_A},\,y\in\mathcal A^{n_B}}
       \prod_{i,j}\beta(x_i,y_j)^{B_{ij}}.
\]
Sum first over \(x\). Character orthogonality gives \(m^{n_A}\)
precisely for \(By=0\) and zero otherwise. This proves the first
formula, including its phase.

For the second, \(4\) acts invertibly on \(\mathcal A\), so the
integral homology theorem gives
\[
 |\Hom(H_1,\mathcal A)|
 =m^{b_G}|\ker A_\epsilon|.
\]
The latter kernel is \(\ker B\oplus\ker B^T\). Perfect pairings
identify the Pontryagin dual of \(\coker B\) with \(\ker B^T\),
and finite cardinalities give
\(|\ker B^T|=m^{n_A-n_B}|\ker B|\).
Substitution, using \(b_G=|E|-n_A-n_B+1\), yields the square of
the positive first formula.
\end{proof}

\section{A concrete strengthened collision}
\begin{corollary}[An eight-vertex pair]
On \(K_{4,4}\), use \(D_0=I_3\) and \(D_1=J_3-I_3\), where
\(J_s\) denotes the all-ones square matrix. The resulting two
oriented irreducible manifolds have
\begin{align*}
 H_1(M_0;\mathbb Z)&\cong\mathbb Z^9\oplus(\mathbb Z/4)^2
                                  \oplus(\mathbb Z/8)^6,\\
 H_1(M_1;\mathbb Z)&\cong\mathbb Z^9\oplus(\mathbb Z/4)^2
                  \oplus(\mathbb Z/8)^4\oplus(\mathbb Z/16)^2.
\end{align*}
They have equal oriented \(\mathrm{SU}(2)\) RT invariants at every
level, including all block-interior colored insertions described above.
They also have equal full complex vacuum-divided RT scalars for every
odd-order pointed modular category:
\[
 \mathcal P_{\mathcal A,Q}(M_0)
 =\mathcal P_{\mathcal A,Q}(M_1)=|\mathcal A|^{9/2}.
\]
\end{corollary}
\begin{proof}
The Smith factors of \(D_0\) are \(1,1,1\), those of \(D_1\) are
\(1,1,2\); direct minors or subtraction proves this. The integral
formula gives the displayed groups, hence nonhomeomorphism. Every
torsion factor is a power of two, so a homomorphism to an odd-order
group kills all torsion. The preceding theorems give all quantum
equalities. Irreducibility follows by the same incompressible-Seifert-
piece argument as in Paper 50B; all vertex degrees here are four.
\end{proof}

\paragraph{A precise detecting experiment in a different TQFT.}
The standard untwisted finite gauge theory~\cite{dw} for a finite group \(K\)
has connected closed-manifold partition function
\(Z_K(M)=|\Hom(\pi_1M,K)|/|K|\); its vacuum-divided version is the
numerator. For \(K=\mathbb Z/2^a\), write that numerator as
\(\mathcal D_a(M)\). For the pair just constructed,
\[
 \mathcal D_a(M_0)=\mathcal D_a(M_1)\quad(a=1,2,3),\qquad
 \mathcal D_4(M_0)=2^{58},\quad\mathcal D_4(M_1)=2^{60}.
\]
Thus the \(\mathbb Z/16\) gauge theory distinguishes them by the
factor four; \(\mathbb Z/2,\mathbb Z/4,\mathbb Z/8\) do not.
This uses a known finite-gauge-theory observable, not a newly claimed
formula for such theories.

\begin{corollary}[Arbitrarily deep abelian probe agreement]
For every \(L\ge1\), one fixed complete bipartite backbone admits
two nonhomeomorphic manifolds which have the same entire
\(\mathrm{SU}(2)\) sequence, the same RT scalar for every odd-order
pointed category, and the same vacuum-divided untwisted
\(\mathbb Z/q\) gauge invariant whenever \(v_2(q)\le L+3\).
They are first distinguished within the \(\mathbb Z/2^a\) hierarchy
at \(a=L+4\), by a factor four.
\end{corollary}
\begin{proof}
Set \(N=2^{L+1}+2\) and, with square sizes indicated, choose
\[
 D_0=(J_{2^L+1}-I_{2^L+1})\oplus I_{2^L},\qquad
 D_1=J_{2^{L+1}+1}-I_{2^{L+1}+1}.
\]
The matrix \(J_s-I_s\) has Smith factors
\(1,\ldots,1,s-1\). Hence both homology groups have the common
part
\(\mathbb Z^{(N-1)^2}\oplus(\mathbb Z/4)^2
\oplus(\mathbb Z/8)^{2(N-2)}\); their last two cyclic factors
are respectively \((\mathbb Z/2^{L+3})^2\) and
\((\mathbb Z/2^{L+4})^2\). All assertions now follow from
\(|\Hom(\mathbb Z/2^e,\mathbb Z/q)|=2^{\min(e,v_2(q))}\).
\end{proof}

\section{Nonabelian odd nilpotent probes}
The preceding examples satisfy \(\det A_\epsilon=\pm2^s\). This
condition forces a stronger invisibility than odd abelian homology alone.
Let \(\pi:\pi_1(M_{G,\epsilon})\twoheadrightarrow F_{b_G}\) be the
canonical graph quotient, which kills every vertex group in the
graph-of-spaces presentation.

\begin{theorem}[Coprime odd nilpotent representations factor through the graph]
Suppose the cap orders are two and \(\det A_\epsilon\ne0\).
For every finite nilpotent group \(K\) of odd order such that
\(\gcd(|K|,\det A_\epsilon)=1\),
precomposition with \(\pi\) gives a bijection
\[
 \Hom(F_{b_G},K)\xrightarrow{\ \cong\ }
 \Hom(\pi_1(M_{G,\epsilon}),K).
\]
Consequently, for every \(\omega\in H^3(K;U(1))\), the untwisted or
twisted Dijkgraaf--Witten invariant is
\[
 Z_{K,\omega}(M_{G,\epsilon})=|K|^{b_G-1},\qquad
 \frac{Z_{K,\omega}(M_{G,\epsilon})}{Z_{K,\omega}(S^3)}=|K|^{b_G}.
\]
For every odd prime \(\ell\nmid\det A_\epsilon\), the pro-\(\ell\) completion of
\(\pi_1(M_{G,\epsilon})\) is the free pro-\(\ell\) group of rank
\(b_G\). In particular, if \(\det A_\epsilon=\pm2^s\), all of the
finite-group assertions hold for every finite odd-order nilpotent
\(K\), and the completion assertion holds for every odd prime.
\end{theorem}
\begin{proof}
Let a homomorphism to \(K\) be fixed. Use lower-case symbols for its
images. Since the exponent of \(K\) is odd, squaring is a bijection:
if \(2u\equiv1\) modulo that exponent, the unique square root of
\(h\) is \(h^u\). The cap relations give
\(a_v=h_v^{-u}\), \(b_v=h_v^u\), hence \(a_vb_v=1\).
The vertex boundary relation therefore becomes
\[
                  \prod_{e\ni v}z_{ve}=1.
\]
In a graph-of-groups presentation after choosing a maximal tree,
each \(z_{ve}\) is a conjugate of \(h_w^{\epsilon_e}\); the
conjugating element is either trivial or a cycle stable letter.

Write \(\gamma_1K=K\) and
\(\gamma_{j+1}K=[K,\gamma_jK]\). Suppose inductively that every
\(h_v\) lies in \(\gamma_jK\), as is automatic for \(j=1\).
In the finite odd abelian quotient \(\gamma_jK/\gamma_{j+1}K\),
conjugation by every element of \(K\) is trivial. The vertex product
relations thus reduce exactly to
\[
                 A_\epsilon\,\overline{\mathbf h}=0.
\]
Multiplication by the adjugate matrix implies
\((\det A_\epsilon)\overline{\mathbf h}=0\). The coprimality
assumption makes multiplication by \(\det A_\epsilon\) invertible
on this quotient, so every \(h_v\) lies in
\(\gamma_{j+1}K\). Nilpotence terminates this induction with
\(h_v=1\). The cap and seam relations now kill \(a_v,b_v,z_{ve}\),
so every vertex-group generator has trivial image. Only the
\(b_G\) stable letters remain, with no relations among them.
This proves the factorization and bijection, including its uniqueness.

Collapsing the capped block interiors to vertices and their seam
collars to intervals realizes the graph quotient by a continuous
map from the manifold to the backbone graph.
Every corresponding classifying map factors, up to homotopy, through
\(BF_{b_G}\), a graph. Its pullback of every degree-three class
\(\omega\) is therefore zero, and the Dijkgraaf--Witten action has
weight one. There are \(|K|^{b_G}\) homomorphisms; the standard
groupoid normalization divides by \(|K|\). Finally the bijection
for all finite \(\ell\)-groups identifies the corresponding pro-
\(\ell\) completions, proving the last assertion.
\end{proof}

Thus the eight-vertex pair, and the arbitrarily deep families above,
also collide for every twisted finite gauge theory whose gauge group is
odd and nilpotent, including nonabelian odd prime-power groups. This
assertion does not extend here to all odd-order groups: removing
nilpotence would invalidate the terminating lower-central-series step.

\paragraph{Scope.}
These are exact extensions of the same construction. They show that
adding all odd pointed theories, or arbitrarily many block-local
\(\mathrm{SU}(2)\) Wilson insertions, need not remove its blindness.
They do not establish equality for every modular category, every
cross-seam colored ribbon, arbitrary finite gauge groups, or linking
pairings. No numerical check is used as a proof.

